% !TeX root = ../papers/nc-p3.tex
\section{反転を伴う固有値衝突の全流入ファイバー}
\label{sec:reflected-collision-fibers}
\begin{theorem}\label{thm:reflected-collision-fibers}
定理\ref{thm:split-repeated-surface}第二行の閉軌道を用い、$t^4\ne1$, $t c\ne0$,
$f=(t^2-1)/(2t)$, $u=f^2$, $v=fc$ とする。
$c^2\ne1$ ならその上の全窓の添字保持同型類は、$u\ne1$ でちょうど五つ、
$u=1$ でちょうど三つである。正規表示の台は、二座標軸について各々
正負のどちらかまたは空を選んだものとする。一斉符号反転で割り、$u=1$ ではさらに
平面の四分の一回転で割る。特殊化は台を小さくすることとこの同値で尽きる。
$c^2\ne-1$ なら全て最小周期2である。$c^2=-1$ なら両軸を使って $u\ne1$ のときだけ
最小周期2、それ以外は1である。

$c^2=1$ では次の16個の相異なる指標を各一次元とする。
\[
 H=\{(\pm2,0),(0,\pm2),(\pm2,\pm2),(\pm3,\pm1),(\pm1,\pm3)\}.
\]
許す非零台は $H$ のうちある開半平面に入る任意の部分集合である。
これらの座標空間を指標通りの二次元トーラスと下記の明示的な有限作用で割ったものが
全無印同型類である。原点を含む全窓が十五次元成分に属し、標準Koszul左右AS正則である。
極大の座標空間は八次元が16個で、有限作用を除くと $u\ne1$ で四つ、$u=1$ で二つ、
各々の一般同型類の次元は6である。各軌道の半安定特殊化は有限作用と台の指標錐の面で決まる。
最小周期は下記の符号表と有限同型検査で2または4に決まり、各極大台の一般点では4である。
\end{theorem}
\begin{proof}[全ての横断指標と非例外値]
位置2,3でビットを反転した既出の28語表示を使う。語 $q$ の重みは
$h=(n_0-n_2,n_1-n_3)$ で、$n_j$ は文字 $j$ の個数である。
全基底群のアフィン軌道次元は59、自明指標の軌道次元は11である。
二切断の接核をこの軌道接で割ると、自明指標が常に二つあり、非自明な指標は
$c^2\ne1$ で四軸指標 $(\pm2,0),(0,\pm2)$ 各一つ、$c^2=1$ で上記 $H$ の各一つとなる。
\path{computations/verify_reflected_collision_fibers.py} は各指標で二切断の左右核の全積を行とし、
分数体での階数を計算する。一方 $tc(t^4-1)$ を逆にした全域で軌道行列の必要小行列式の
同時零点がないことを確認する。接行列はさらに $c^2-1$ を逆にした場合と
$c=1$ に制限した場合で同じ検査を行う。一移動で $c=-1$ に移せる。
下界検査は前節と同じ単元ピボット消去および小行列式イデアルの局所化による。
従って全ての階数が各指定範囲で一定である。アフィン安定群は自明スカラー三次元と
上記トーラスで恒等成分を尽くす。

非例外値では任意の正台を含む余指標は $(\pm1,\pm1)$ のいずれかでよい。
全正座標は余重み2で、最大語重みは4なので二次補正だけを要する。
$M$ の一般逆を $g$、左右核を $L,R$ とすると二座標の代表 $A,B$ の補正右辺は
$L(AgB+BgA)R$ である。検算は $\mathbb Q(t,c)$ で各軌道余空間の一本を選び、
両軸の平方および二種類の象限の混合積の全四列が対応する接行列の像に入ることを確かめる。
反対象限は全符号反転で移る。定階数の接核・余核上での二次障害は一般体上で零なので全域で零である。
余重み4に新たな横断接はなく、各象限で二成分とも自由である。トーラスは非零な二値を独立に規格化する。
\end{proof}
\begin{proof}[例外値で全ての正方向が積分できる理由]
全階数スキームの軌道横断面を取り、接射影の座標で滑らかな18変数形式空間に埋め込む。
二つの中性座標は $t,c$ で張れる。実際、28語上の既出の対角不変比 $u,v$ の微分は
全域で独立であり、軌道の重み零部分は対角作用である。
既約なアフィン76次元の十五次元成分は $c=1$ の閉軌道を通るので、
その被約横断面は純次元17である。形式的な軌道との積分解と有限型代数の完備化の
被約性・等次元性による。従って正則形式局所環でのこの横断面のイデアルは高さ1の
素イデアルの交わりであり、一つの非零元で生成される。
トーラス同変な生成元を選べる。実際、各有限jetで一次元の極小生成空間の
固有ベクトルを同じ重みで両立して持ち上げ、中山の補題を使えばよい。
この生成元の重みは零である。さもなくば中性二平面への制限が恒等零となる。
中性平面には上記二変数の固定点族が形式同型に入り、$c^2\ne1$ の一般点は
定理\ref{thm:split-repeated-surface}によりこの成分の外にあるので、これは不可能である。

従って中性座標を原点に固定し、一余指標で正になる座標のみを任意に取れば、
この重み零の方程式は定数項以外を持てず恒等零となる。これで全ての正座標が同時に積分できる。
実際のテンソルは定理\ref{thm:mixed-flip-fiber}のSchur再帰と同様、
軌道接の補および接核の補を同変に選び、正座標の単項式ごとに
$L(D_m-\sum_{a+b=m}D_a g B_b)R=0$ を解いて得る。
正余重みは有界なのでこれは多項式であり、入力の接射影は指定した座標そのものとなる。
Luna横断面での流入弧の一意性と閉軌道への一径数判定で全体を覆い、
重なりでは同じ台の同じ点になる。全次数性は正則な閉点への良い還元から従う。
\end{proof}
\begin{proof}[有限比較と周期]
四相異なる指標は全位置の四文字を区別し、恒等安定群の正規化群は各位置で単項である。
支持にない隣接対は同一文字の対だけなので各位置の置換は同じである。
$u,v$ の比較により、ある固定基点の置換は一般に $-1$（二軸の同時反転）までで、
$c^2=1$ なら軸交換、$u=1$ なら四分の一回転も許す。
これは前の八置換をこの二不変量に作用させた結果である。$c^2\ne1$ での個数が従う。
この場合二移動は置換なしで可能で、一移動は $c^2=-1$ でちょうど一軸の符号反転として実現する。
台の二つの位相はトーラスで自由に変わるので、台だけで判定して表示通りの周期となる。

例外値での有限作用を係数まで指定する。$N_{q,ij}=[q_i=j]$ を28語の行と16文字の列で作ると、
その非零Smith因子は11個全て1である。許した置換 $p$ を実装するには
比 $w_{p q}/w_q$ の列に、$N$ のSmith左変換の整数行を指数とする冪積を取る。
最初の11積を残し、残る対角パラメータを1とし、Smith右変換の整数行の冪積で16個の係数 $d_{i,j}$ に戻す。
この単項作用は語 $q$ を $p q$ へ移し $\prod_i d_{i,q_i}$ 倍する。他の解との差は恒等トーラスだけである。
各 $h$ の座標線の基底は、接行列の核の通常の自由列基底で軌道列の張る空間の外にある最初の列の剰余類
とすればよい。上の単項作用をその列に掛け、行消去で新しい重みでの軌道列を引いた一係数が
座標間の倍率である。横断面と接核への射影を同変に取るので、この剰余作用だけで全座標に作用する。
トーラスについては、同じ台での係数比の冪積を指標行列の整数核基底上で比較すればよい。
従って有限候補と線型演算・整数冪積のみの同型判定であり、二つの基点の表示が異なる場合は
前の $u,v$ の基点同型を一つ作用させて同じ判定を行う。

$H$ の八つの反対対から壁の16個の室を得る。各室にちょうど八指標が正である。
$u\ne1$ の四元群、$u=1$ の二面体群はいずれも室へ自由に作用する。
各室の全台の軌道次元は2なので表示した極大族と次元を得る。
トーラス軌道の閉包の面判定と有限比較から軌道の特殊化も従う。

$c=1$ の二移動は、語を二位置巡回した後、旧位置0,1の文字0,1に $f$、
旧位置2,3の文字0,1に $1/f$ を掛けて補正する。これは背景を保つ対合である。
接座標で負号となる $h$ は
\[
       (-3,1),\ (-2,0),\ (-1,3),\ (1,3),\ (2,0),\ (3,1)
\]
だけで、他は正号である。各指標の接商は全域で直線束なので、対合の符号は一定である。
上記検算は $t=2$ で軌道列と核基底からこの全符号を有理数上で求める。
従って座標にこの六符号を掛け、前段の有限同型検査を通れば最小周期2、他は4である。
$c=-1$ では一移動で $c=1$ へ移して使える。極大台を保つ有限置換は恒等だけで、
二つの軸での符号は $-,+$、それらの和または差の指標では全て $+$ なのでトーラス倍率では一致しない。
これで全台での最小周期と一般周期が決まる。
\end{proof}
